\subsection{Vacua with Small Flux Superpotential --- arXiv:1912.10047}

\textbf{Authors:} Mehmet Demirtas, Manki Kim, Liam McAllister, Jakob Moritz

\paragraph{主な主張}
Type IIBのCalabi--Yau orientifoldにおいて、大複素構造・弱結合領域で指数関数的に小さいフラックス超ポテンシャルを構成し、\(\lvert W_0\rvert\simeq 2\times10^{-8}\)の具体例を示した。 \cite{Demirtas:2019sip}

\paragraph{新規性}
摂動的に平坦で超ポテンシャルが消えるフラックス配置を先に構成し、periodの指数関数的補正で平坦方向を固定することにより、小さい\(W_0\)を系統的に探索可能にした。

\paragraph{位置付け}
KKLT構成\cite{Kachru:2003aw}の小さい\(W_0\)という入力を具体化する研究であり、全K\"ahlerモジュライの固定やde Sitter uplift自体は今後の課題として残す。

\paragraph{コメント（読書メモ）}
檜垣さんの科研費申請書の参考資料[4]。元メモでは「弦と膜が生み出す量子効果を考慮すると、小さなポテンシャルが生成し得る」と整理した。Gukov--Vafa--Witten超ポテンシャルの真空値を十分小さくする模型で、\(\lvert W_0\rvert\ll1\)は制御されたKKLT構成に必要な条件として読む。
